%~ \documentclass[12pt]{article} 
%~ \usepackage{bm} 
%~ \usepackage{graphics,graphicx} 
%~ \usepackage{pstricks} 
%~ \usepackage{amssymb,amsmath} 
%\usepackage{color} \pagestyle{empty} 
%Paper size is 8 x 11 1/2 inches. One inch margin at the top and the 
%right. 
%~ \setlength{\oddsidemargin}{-0.25in} 
%~ \setlength{\topmargin}{0in} 
%~ \textheight 9.0in 
%~ \textwidth 6.0in 
%~ \renewcommand{\baselinestretch}{1.0} 
%~ \addtolength{\topmargin}{-0.5in} 
%~ \def\by#1#2{{\displaystyle {#1}\over \displaystyle {#2}}} 
%~ \def\d{{\rm d}} 

%\begin{document}
%\noindent\includegraphics[width=1.9cm]{logo.ps}\hfill
\begin{center}
{\LARGE  \sf {The Institute of Mathematical Sciences}}\\
\rule{\textwidth}{1pt}
(\sf{An autonomous institute under the }\\
\sf{Department of Atomic Energy})\\
\sf{C.I.T.  Campus, Taramani PO, Madras 600 113, India}
\end{center}

\begin{flushleft}
{\large\sf Dr. G. Rajasekaran}\hfill {\sf 
\today}\\
\vskip 0.5cm
{\large\sf Dr. M.V.N. Murthy} 
\end{flushleft}

\hrule 
%The address is put in an add.tex file
\begin{flushleft}
Dr. M.R. Srinivasan\\
Member, AEC\\
"SUNNINGDALE", Kotagiri Road\\
OOTACAMUND - 643 002\\
The Nilgiris
\end{flushleft}
\parskip 8pt
\vskip .3cm
%input matter

\noindent Dear Dr. Srinivasan,

As you are well aware, the India-based Neutrino Observatory (INO) is a
mega science flag ship project of the Department of Atomic Energy and is
also supported by the Department of Science and Technology. In the 2012
Indian Science Congress, the Prime Minister Dr Manmohan Singh announced
that "the Government is also considering a proposal to establish a
Neutrino Observatory in Theni district in Tamil Nadu with a proposed
investment of Rs.1350 crore."

The INO project has multiple goals: to study atmospheric neutrinos,
neutrinoless double beta decay, dark matter and low energy neutrino
spectroscopy. Each one of these goals will add to the knowledge at the
cutting edge of high energy physics and particle astrophysics.

In particular, the R\& D for the magnetised Iron Calorimetric detector
(ICAL), which is designed to study the properties of atmospheric neutrinos
in the first phase, is complete and steps towards scaling up to the
actual size with industrial interface are being taken. When completed
this will be the largest magnetised detector in the world. The main goal
of this large detector is the precision study of neutrino parameters
which are not yet well understood.

Of particular fundamental importance is the determination of the ordering
of the neutrino masses or what is known as the hierarchy problem. Until
last year it was not clear that this could be done at ICAL within a
reasonable time of few years. However, the measurement of a parameter,
called theta-13, by the Chinese, Korean and French groups has brought ICAL
to the forefront of the future neutrino experiments. It is now recognised
worldwide that ICAL will be a unique facility to solve the problem of
hierarchy, buttressing its discovery potential. More generally, INO offers
a unique window towards physics beyond the Standard Model of particle
physics which got a resounding confirmation through the recent discovery
of the Higgs boson.

It is now clear that worldwide there are two proposed experiments which
can determine the neutrino mass hierarchy in a completely self-contained
manner: the ICAL at INO and the proposed Daya Bay-2 in China. The Daya
Bay-2 proposal is expected to take data from 2019 and would take about
five years to claim the results for hierarchy with reasonable confidence
assuming that there are no unexpected delays.

In order for ICAL at INO to be relevant and competitive, it is necessary
that the detector is installed before 2019 and start taking data. This
can be achieved if the making of the tunnel and cavern and outfitting
are completed before the end of the year 2017.

The INO collaboration was informally formed in 2001 and formalised in 2003
through an MOU. Beginning with a small number of scientists it is now
about 100 strong. The members of the collaboration have worked very hard
to bring INO to its present level. The collaboration includes nearly 20
young students who are part of the graduate training program initiated in
2009 to train young scientists adept at experimental and theoretical
aspects of neutrino physics. To our knowledge, no other basic science
project has produced such an extensive body of work or human resource
before making a plea for its approval and execution. It is also
for the first time that such a large group of scientists have come
together from so many different institutions in India to work towards
a common goal.

Therefore any further delay and dilution of the objectives will only
create more frustration and there is now a real danger that if the project
is not approved immediately in the next few months, the collaboration may
in fact collapse.

Hence it is imperative that the funding agencies and authorities in charge
repose confidence in the INO collaboration and take steps towards a speedy
implementation of the project.

There is thus a huge challenge and an opportunity here to demonstrate that
a large experiment done in India can indeed succeed at the cutting edge of
science. The immediate tasks are:

* To establish the INO with urgency and create a structure that can take
quick decisions with enough financial strength to execute the engineering
and scientific challenges ahead for construction of the laboratory and the
detector in a timely manner that does not compromise its objectives.

* Begin the construction of the underground laboratory at the earliest-
the engineering DPR clearly lays down the time line of four years- if
promptly executed the laboratory and its infrastructure can be ready
before the end of 2017.

* In parallel, the various components of the ICAL detector: iron plates,
RPC detectors, coils, and electronics may be assembled for in situ
construction of the detector once the lab is ready. The R\& D phase has
been completed and much work on scaling up, with industrial involvement,
has already been done.

* Recruitment of scientists and engineers to complete this job must begin
immediately. The objective cannot be achieved without human resources.
The management of INO should have the freedom to recruit people directly
if this objective is to be achieved.

We hope you will use your position and authority to enable the goals of
this ambitious experiment to be realised.

Thank you very much. Best regards,

\begin{flushleft}
Yours sincerely\\[1.5cm]
G Rajasekaran\\
Professor Emeritus\\
The Institute of Mathematical Sciences\\
Chennai 600113\\
Adjunct Professor\\
Chennai Mathematical Institute\\
Siruseri.\\
\vskip 1cm 
M V N Murthy\\
Professor\\
The Institute of Mathematical Sciences\\
Chennai 600 113.\\
\end{flushleft}

\vfill

\hrule
\begin{flushleft}
{\sl Phone:~(91 44)2254 3326,\hfill Fax:~(91 44)2254 1586}\\[5pt]
{\sl Email:~graj@imsc.res.in} 
\end{flushleft}

%\end{document}


